% Paste from here into your main Overleaf document.
% Required packages (add to your preamble if not already there):
%   \usepackage{booktabs}
%   \usepackage{enumitem}

%%-----------------------------------------------------------------------
\chapter{EGON noter}
\section{Strukturovervejelser}

\subsection{Hvad Thomas sandsynligvis mente}

\begin{itemize}
    \item \textbf{ARIMA} $\rightarrow$ baseline på dansk inflation
    \item \textbf{GETS $\rightarrow$ Dynamic regression eller VAR} med dansk arbejdsløshed på dansk inflation. Det er den klassiske Phillips curve-test: har arbejdsløshed forecasting-værdi for inflation?
    \item \textbf{Granger-kausalitet med EU-data} som et \textit{analytisk element} -- altså ikke en tredje model, men en test der besvarer: er dansk inflation påvirket af europæiske forhold? Det bruges til at \textit{motivere} eller \textit{diskutere} modellen, ikke som en selvstændig forecasting-model.
\end{itemize}

\subsection{Hvad det betyder for den tidligere anbefaling}

Den tidligere anbefaling var ikke så langt fra, men EU-delen blev overkompliceret ved at smide Eurozone-inflation ind i selve VAR-modellen. Det er nok for meget.

En renere og skarpere struktur ville være:

\begin{enumerate}
    \item \textbf{ARIMA} -- univariat baseline på dansk inflation
    \item \textbf{GETS $\rightarrow$ Dynamic regression eller VAR} -- dansk inflation og dansk arbejdsløshed (Phillips curve)
    \item \textbf{Granger-kausalitet} som et diskussionselement: påvirker EU-inflation dansk inflation? Det kan bruges i data-afsnittet eller i diskussionen til at kontekstualisere resultaterne
\end{enumerate}

Det er en mere fokuseret opgave, og man undgår at drukne i for mange variable.

\section{Research Question}

\textbf{Does the Phillips Curve have forecasting value for Danish inflation,
and does incorporating European data improve forecasts relative to a
univariate benchmark?}

The Phillips Curve describes the empirical relationship between inflation
and unemployment. From a forecasting perspective it raises a concrete and
testable question: can a model that also uses unemployment outperform a
pure time-series model of inflation? This is the central question of the
assignment.

The Denmark versus Eurozone angle sharpens the analysis further. Denmark
is a small, open economy tightly integrated with the European business
cycle, so European inflation and unemployment may Granger-cause Danish
inflation. A bivariate or multivariate model that incorporates EU data
may therefore forecast better than a model based on domestic data alone.
At the same time, Denmark's ``flexicurity'' labour market institutions
may produce a structurally different unemployment-inflation tradeoff from
the Eurozone average, giving us something meaningful to compare and
discuss.

The relationship is also known to be \textbf{unstable over time}. The
Phillips Curve ``disappeared'' in many countries during the 1980s and
1990s, and there is ongoing debate about whether it has re-emerged
post-COVID. Formal tests for structural breaks are therefore not an
add-on but a core part of the analysis. Key references are
\textcite{stockwatson1999} and \textcite{galigertler1999}.

%%-----------------------------------------------------------------------
\section{Data}

\subsection{Series}

Going back to 1975 rather than the early 1990s significantly strengthens
the analysis. It adds the oil shocks of 1973 and 1979, the Volcker
disinflation of the early 1980s, and the Danish ERM exit of 1992 as
additional structural break candidates, and gives the Bai-Perron test
roughly 200 quarterly observations to work with.

\begin{table}[h!]
\centering
\begin{tabular}{llll}
\toprule
\textbf{Series} & \textbf{Frequency} & \textbf{Period} & \textbf{Source} \\
\midrule
Danish unemployment rate   & Quarterly & 1975 onward & Statistics Denmark (dst.dk) \\
Danish HICP / CPI          & Quarterly & 1975 onward & Statistics Denmark / Eurostat \\
Eurozone unemployment rate & Quarterly & 1975 onward & ECB AWM database / OECD \\
Eurozone HICP              & Quarterly & 1975 onward & ECB AWM database / Eurostat \\
\bottomrule
\end{tabular}
\caption{Data overview. The ECB Area-Wide Model (AWM) database provides
a synthetic Eurozone aggregate constructed back to 1970, enabling a
consistent long-run comparison.}
\end{table}

\subsection{Transformations}

\begin{itemize}[nosep]
    \item Inflation is computed as the year-on-year percentage change in
          the price index.
    \item Unemployment is used as a rate in levels; stationarity is
          assessed empirically via ADF and KPSS tests.
    \item First differences are applied to any series found to be
          non-stationary.
    \item No population or calendar adjustments are expected to be
          necessary, but this is verified during exploratory analysis.
\end{itemize}

%%-----------------------------------------------------------------------
\section{Methods}

The core forecasting question is whether a model that incorporates
unemployment can outperform a univariate ARIMA benchmark. The analysis
builds toward that comparison through the following stages.

\subsection{Stage 1: Exploratory Analysis}

\begin{itemize}[nosep]
    \item Time plots for all four series, with visual inspection for
          trend, structural shifts, and unusual observations.
    \item STL decomposition to assess seasonal patterns in unemployment.
    \item ACF and PACF plots to characterise autocorrelation structure.
    \item ADF and KPSS unit root tests to determine the order of
          integration for each series.
\end{itemize}

\subsection{Stage 2: Structural Break Tests}

The instability of the Phillips Curve is well-documented in the
literature and is a central theme of the paper. We test for it formally
rather than assuming it away.

\begin{itemize}[nosep]
    \item \textbf{Chow test} at candidate break dates: the first oil
          shock (1973 Q4), the Volcker disinflation (1982 Q1), the Danish
          ERM exit (1992 Q3), the global financial crisis (2008 Q3), the
          COVID shock (2020 Q1), and the post-pandemic inflation surge
          (2021 Q3).
    \item \textbf{Bai-Perron multiple breakpoint test} to let the data
          determine the number and timing of breaks endogenously, without
          imposing candidate dates.
    \item \textbf{CUSUM and CUSUM-of-squares tests} to detect gradual
          parameter instability over the full sample.
    \item Results directly inform model estimation: we either include
          break dummies, split the sample, or restrict the estimation
          window to the most recent stable regime.
\end{itemize}

\subsection{Stage 3: General-to-Specific Model Selection (GETS)}

GETS is particularly well-suited here because the candidate regressor
set is wide: Danish inflation, Eurozone inflation, Danish unemployment,
Eurozone unemployment, and lagged values of each. Starting broad and
testing down finds the parsimonious model the data actually supports,
rather than imposing a specification in advance. This is the approach
advocated by \textcite{doornik2018} and is one of the methodological
elements the lecturer specifically highlighted as underused.

\begin{itemize}[nosep]
    \item Specify a General Unrestricted Model (GUM):
          \[
          \pi_t^{DK} = \alpha
            + \sum_{i=1}^{p} \beta_i \pi_{t-i}^{DK}
            + \sum_{i=1}^{p} \delta_i \pi_{t-i}^{EU}
            + \sum_{j=0}^{q} \gamma_j u_{t-j}^{DK}
            + \sum_{j=0}^{q} \phi_j u_{t-j}^{EU}
            + \varepsilon_t
          \]
          with initial lag length set generously ($p = q = 4$).
    \item Apply sequential $t$-tests and $F$-tests to eliminate
          insignificant regressors, verifying residual well-behavedness
          at each step.
    \item AIC and BIC serve as a cross-check on the final specification.
    \item Run the same procedure for Eurozone inflation to compare
          resulting model structures across the two regions.
\end{itemize}

\subsection{Stage 4: Benchmark Models}

\begin{itemize}[nosep]
    \item Na\"{i}ve forecast.
    \item Seasonal na\"{i}ve forecast.
    \item Simple exponential smoothing and Holt's method.
\end{itemize}

\subsection{Stage 5: Univariate ARIMA}

This is the primary benchmark that the multivariate models must beat
to demonstrate that unemployment contains forecasting information
beyond what is already in past inflation.

\begin{itemize}[nosep]
    \item Fit ARIMA$(p,d,q)$ to Danish and Eurozone inflation separately,
          using the Box-Jenkins procedure.
    \item Select the final specification via AIC and BIC.
\end{itemize}

\subsection{Stage 6: Dynamic Regression (ARDL)}

\begin{itemize}[nosep]
    \item Estimate the parsimonious model from Stage 3 as a dynamic
          regression with unemployment (domestic and/or European) as
          external regressors.
    \item Incorporate break dummies where Stage 2 identified instability.
    \item Residual diagnostics: Ljung-Box test for autocorrelation, ARCH
          test for heteroskedasticity.
\end{itemize}

\subsection{Stage 7: Vector Autoregression (VAR)}

\begin{itemize}[nosep]
    \item Specify a bivariate VAR in $(\pi_t^{DK},\, u_t^{DK})$ and a
          four-variable VAR in
          $(\pi_t^{DK},\, \pi_t^{EU},\, u_t^{DK},\, u_t^{EU})$.
    \item Select lag length via AIC / BIC / HQIC.
    \item Granger causality tests: does EU inflation Granger-cause Danish
          inflation? Does unemployment Granger-cause inflation?
    \item Impulse response functions to trace how a shock to European
          unemployment propagates to Danish inflation.
    \item Multi-step forecasts from the VAR as an additional competitor
          model.
\end{itemize}

\subsection{Stage 8: Forecast Evaluation}

\begin{itemize}[nosep]
    \item Hold out the final 8 quarters as an evaluation period.
    \item Generate one-step-ahead and four-step-ahead forecasts from all
          models using a rolling window.
    \item Compare accuracy via RMSE and MAE.
    \item \textbf{Diebold-Mariano test} to assess whether the best
          multivariate model statistically outperforms the ARIMA
          benchmark. This is the direct answer to the research question.
    \item Report 95\% prediction intervals and assess empirical coverage.
\end{itemize}

%%-----------------------------------------------------------------------
\section{Paper Structure}

\begin{enumerate}[nosep]
    \item \textbf{Introduction} Research question, motivation, and
          overview.
    \item \textbf{Background} Phillips Curve theory, evidence on
          instability, Denmark versus Eurozone labour market
          institutions.
    \item \textbf{Data} Sources, construction of long-run series,
          descriptive statistics.
    \item \textbf{Exploratory Analysis} Time plots, decomposition,
          ACF/PACF, unit root tests.
    \item \textbf{Structural Break Analysis} Methodology and results;
          implications for estimation.
    \item \textbf{Model Selection} GETS procedure and final
          specifications for both regions.
    \item \textbf{Model Estimation} ARIMA, dynamic regression, and VAR
          results with diagnostics.
    \item \textbf{Forecast Evaluation} Out-of-sample comparison,
          Diebold-Mariano test, prediction intervals.
    \item \textbf{Discussion} Does unemployment help forecast inflation?
          What do the cross-country differences reveal? Limitations.
    \item \textbf{Conclusion} Main findings and directions for further
          work.
\end{enumerate}

%%-----------------------------------------------------------------------
\section{Division of Work}

\begin{table}[h!]
\centering
\begin{tabular}{ll}
\toprule
\textbf{Section} & \textbf{Responsible} \\
\midrule
Data collection and cleaning         & TBD \\
Exploratory analysis (Stage 1)       & TBD \\
Structural break tests (Stage 2)     & TBD \\
GETS model selection (Stage 3)       & TBD \\
ARIMA and benchmarks (Stages 4 and 5) & TBD \\
Dynamic regression (Stage 6)         & TBD \\
VAR (Stage 7)                        & TBD \\
Forecast evaluation (Stage 8)        & TBD \\
Writing and editing                  & All \\
\bottomrule
\end{tabular}
\caption{Work distribution (to be finalised by the group)}
\end{table}

%%-----------------------------------------------------------------------
\section{Software and Replication}

All analysis will be conducted in \textbf{R}. Key packages:
\texttt{forecast} and \texttt{fable} (ARIMA, exponential smoothing,
forecast evaluation), \texttt{strucchange} (Chow, Bai-Perron, CUSUM),
\texttt{gets} (general-to-specific selection), \texttt{vars} (VAR,
Granger causality, IRFs), and \texttt{dynlm} (dynamic regression).
Code will be shared via a joint GitHub repository to ensure full
replicability.

% End of paste